\documentclass[10pt,a4paper]{amsart}
\usepackage[margin=19mm]{geometry}
\usepackage{amsmath,amssymb,mathtools}
\usepackage[scr=rsfs,cal=euler]{mathalfa}
\usepackage{xfrac}
\usepackage[unicode,hidelinks,bookmarks=false]{hyperref}
\newcommand{\ie}{i.e.\ }
\newcommand{\cf}{cf.\ }
\newcommand{\resp}{respectively\ }
\newcommand{\Subsection}{\subsection}
\providecommand{\nobreakdash}{\nobreak\hbox{-}}
\setlength{\emergencystretch}{2em}
\multlinegap0pt
\usepackage{bm}
\usepackage{bbm}
\usepackage{dsfont}
\usepackage{xspace}
\usepackage{cancel}
\usepackage{mathtools}
\usepackage{amsthm}
\makeatletter
\@ifundefined{theorem}{\newtheorem{theorem}{Theorem}}{}
\@ifundefined{lemma}{\newtheorem{lemma}[theorem]{Lemma}}{}
\makeatother
\title[Spectral methods for Neural Integral Equations]{Spectral methods for Neural Integral Equations}
\author{Emanuele Zappala}
\date{Source: arXiv:2312.05654v4 (CC BY 4.0)}
\begin{document}
\maketitle

\noindent\emph{Source and licence.} Spectral methods for Neural Integral Equations. Version arXiv:2312.05654v4 (CC BY 4.0), distributed under the Creative Commons Attribution 4.0 International licence, as recorded in the arXiv licence metadata for this exact version and shown on the abstract page https://arxiv.org/abs/2312.05654v4. Full rights notice: licenses/arxiv-2312.05654-CC-BY-4.0.txt.\par\smallskip

\noindent\textit{Editorial scope.} Counted unit: the Theorem whose source label is \texttt{None} in the article (source0.tex lines 412--414, source label \texttt{None}), delivered together with the article's own complete original proof of it (source0.tex lines 415--433, one \texttt{proof} environment of 19 lines). No mathematics is written, repaired or re-derived: statement and proof are the article's own text line for line. Required context retained uncounted: thm:regular\_T (lines 300--310); lem:spec\_approximation (lines 367--373). Every definition, display and label of those lines is retained unchanged. Omitted from this package and not redistributed: 1--299, 311--366, 374--411, 434--1167. Nothing inside a retained excerpt is omitted or altered: every retained body is delivered exactly as the article's lines are written, so no omission receipt of any kind accompanies this package. Citations: the retained excerpts carry no citation command at all, so no bibliography entry of the article is transcribed and no reference is redistributed. Reference adaptation: none was needed and none was made; every reference inside the delivered unit resolves inside it, so no unresolved reference or citation is produced. Printed slips kept literal: no printed word, symbol, hypothesis or number of the counted statement or of its proof was altered. Layout and class: the article's own class is \texttt{amsart} with packages including xcolor, amsthm, amsfonts, babel, graphicx, soul, stfloats, morefloats, cite, lscape, epstopdf, braket, amsrefs, mathbbol; the wrapper is the lane's pinned \texttt{amsart} editorial wrapper at 10pt on A4, which restates the theorem-like environments the retained text needs and adds the article's own macro definitions that the retained text needs (none), with extra wrapper packages (bm, bbm, dsfont, xspace, cancel, mathtools). The delivered numbering is the unit's own. Assets: no retained span contains a figure environment, a graphics-inclusion command, a PDF-page inclusion, a table or a TikZ picture; no image file is copied into this package, the package declares no asset dependency and no image file is redistributed with it. Licence: version arXiv:2312.05654v4 of this article is distributed under the Creative Commons Attribution 4.0 International licence, as recorded in the arXiv licence metadata for this exact version and shown on the abstract page https://arxiv.org/abs/2312.05654v4. A full rights notice is delivered at licenses/arxiv-2312.05654-CC-BY-4.0.txt. Characters: every retained excerpt is pure ASCII, so the delivered engine, XeLaTeX with its default Latin Modern text font, reports no missing character.
	\begin{theorem}\label{thm:regular_T}
		In the hypotheses above, the projected equation 
		$$
		P_NT(y_N) + f = y_N
		$$
		admits at least a solution $u_N$ for all $N$ large enough. Moreover, 
		$$
		\|u - u_N\| \longrightarrow 0
		$$
		as $N \longrightarrow \infty$, with the same rate as $\|P_Nu - u\| \longrightarrow 0$. 
	\end{theorem}

	\begin{lemma}\label{lem:spec_approximation}
		With notation and assumptions as above, let $T$ denote a Fredholm or Volterra integral operator, and let $P_NT : X_N \longrightarrow X_N$ be its projection on $X_N$. Then, for any fixed $\epsilon > 0$, there exist neural networks $F_N, g_N: X_N \longrightarrow X_N$ that satisfiy the condition
		$$
		\|P_NT(y) - P_N \frak I F_N(y) - g_N(y)\|_\infty < \epsilon,
		$$
		for all $y\in X_N$, where $\frak I$ denotes the spectral integration in Fredholm or Volterra form (depending on $T$). Moreover, $P_N \frak I F_N$ approximates with arbitrarily high precision the operator defined through: $P_N\int_{-1}^{\alpha(t)} P_NG(y,t,s)ds$. 
	\end{lemma}

	\begin{theorem}
		Let $T(y) + f = y$ be an integral equation that admits a unique solution $u\in C^{k,\beta}$. Then, for any $\epsilon>0$ we can find $N>0$, and neural networks $G$, $g$ and $\tilde f$ such that if $\tilde u$ is the solution of $P_N\frak I G(y) + g(y) + \tilde f = y$ which we assume to have solutions for $N$ large enough for any $\tilde f$, we have $\|u-\tilde u\|_\infty < \epsilon$.
	\end{theorem}

	\begin{proof}
		For $N$ large enough, $P_N\frak I G(y) + g(y) + \tilde f = y$ has a solution $\tilde u_N$ and we can write 
		$$
		\|u - \tilde u_N\| \leq \|u - u_N\| + \|u_N - \tilde u_N\|,
		$$
		where $u_N$ is the solution to the projected equation $P_NT(y) + P_Nf = y$ as in the case of Theorem~\ref{thm:regular_T}. Then, we can choose $N$ large enough such that $\|u - u_N\|_\infty \leq  \|u - u_N\| < \epsilon/2$, since $\|u-u_N\| \longrightarrow 0$. From Lemma~\ref{lem:spec_approximation} we can find neural networks $G$ and $g$ such that 
		$$
		\|P_NT(y) - P_N \frak I G(y) - g(y)\|_\infty < \epsilon/4,
		$$
		for all $y$ in $X_N$, the projection space of continuous functions (in particular the projection of the H\"{o}lder space). 	Since $f$ is continuous, we can also choose (via the Universal Approximation Theorem) a neural network $\tilde f$ such that
		$$
		\|P_Nf - \tilde f\|_\infty < \epsilon/4. 
		$$
		Then, from the definition of $u_N$ and $\tilde u_N$ we have 
		$$
		\|u_N - \tilde u_N\|_\infty = \|P_NT(u_N) + P_Nf - P_N \frak I G(y) - g(y) - \tilde f\|_\infty \leq \epsilon/4 + \epsilon/4.
		$$
		Therefore, we find that for $N$ large enough, $\|u - \tilde u_N\|_\infty < \epsilon$, which completes the proof. 
	\end{proof}

\end{document}
